\parthead{II}{How a modest bank accumulates a turn}\label{part:2}
In straight, level flight, lift supports the aircraft's weight. Tilt that lift vector by banking the wings and it acquires a horizontal component. The horizontal component bends the velocity; the vertical component must continue to support the aircraft. A turn is the simultaneous performance of those two jobs.

For a steady, coordinated, level turn in still air, write lift as $\mathcal L$, mass as $m$, airspeed as $V$, bank angle as $\phi$ and turn radius as $R$. Resolving forces vertically and towards the centre of the turn gives
\begin{equation}
\mathcal L\cos\phi=mg,\qquad \mathcal L\sin\phi=\frac{mV^2}{R}
\quad\Longrightarrow\quad R=\frac{V^2}{g\tan\phi}.
\label{eq:turn}
\end{equation}
Dividing the two force equations removes both mass and lift. The resulting relation describes the geometry required by speed and bank. An actual aircraft still has to generate the lift: mass continues to matter for lift capacity, stall and structural demands.

\fig{03_turn}{1}{A banked lift vector supports weight vertically and supplies the inward acceleration horizontally. The plan-view arc accumulates $47^\circ$ of direction change at a constant $15^\circ$ bank. At 78 m/s, equation~\eqref{eq:turn} gives a 2.315 km radius, 1.899 km of arc and 24.35 s of turning time, using $g=9.80665$ m/s$^2$. This ideal coordinated, level, still-air benchmark excludes roll-in and roll-out; it is not spatially fitted to the historical approach.}{fig:turn}

Now put 78 m/s and $15^\circ$ into the equation. The radius is approximately 2.32 km. That may initially sound too large for a manoeuvre so often described through images of nearby buildings. But a radius is not the length of the arc that must be flown, and the aircraft is not completing a circle. For a direction change $\Delta\psi=47\pi/180$ radians, the arc length is $R\Delta\psi$, approximately 1,899 m. At the stipulated speed it takes 24.35 seconds.

Fifteen degrees of bank does not produce fifteen degrees of total turn. It produces a turning rate. Here the rate persists long enough to accumulate the much larger change in direction. In the ideal model, $\dot\psi=g\tan\phi/V$; the final angle is its integral over time. The bank angle and the angle turned live in different parts of that calculation.

The load factor offers a second check on intuition. Vertical force balance requires $\mathcal L/(mg)=1/\cos\phi$. At $15^\circ$ it is 1.0353. In this benchmark, the lift magnitude is only about 3.5\% greater than weight. The considerable change in direction comes from sustaining a modest inward acceleration over time, rather than from an instant of extreme force.

Speed alters the bargain. At fixed bank, radius grows as $V^2$: moving faster requires a wider circle. Yet the time for a fixed direction change grows only linearly with $V$, because the greater arc length is traversed at the greater speed. These are two views of the same motion, and both matter to a runway approach. A turn that can be completed quickly may still occupy too much space; a narrow turn may demand more bank than the permitted envelope allows.

\begin{samepage}
An aeroplane also needs time to roll into the bank and out again. During those transitions the turning rate changes. For the complete landing, speed, bank and descent change together, so the circular arc becomes an explanatory benchmark rather than a prescribed route.
\par
\end{samepage}

One requirement is still absent from the diagram. The aircraft has to lose height. It must reduce its downward motion before contact and bring an acceptable forward speed onto the runway. The turn has used almost two kilometres of motion in our example; the next calculation asks what the rest of that motion must accomplish.
